\section{Double complexes and associated total complexes}
\label{section-double-complexes}

\noindent
We discuss double complexes and associated total complexes.

\begin{definition}
\label{definition-double-complex}
Let $\mathcal{A}$ be an additive category.
A {\it double complex} in $\mathcal{A}$ is given
by a system $(\{A^{p, q}, d_1^{p, q}, d_2^{p, q}\}_{p, q\in \mathbf{Z}})$,
where each $A^{p, q}$ is an object of $\mathcal{A}$ and
$d_1^{p, q} : A^{p, q} \to A^{p + 1, q}$ and
$d_2^{p, q} : A^{p, q} \to A^{p, q + 1}$ are morphisms of $\mathcal{A}$
such that the following rules hold:
\begin{enumerate}
\item $d_1^{p + 1, q} \circ d_1^{p, q} = 0$
\item $d_2^{p, q + 1} \circ d_2^{p, q} = 0$
\item $d_1^{p, q + 1} \circ d_2^{p, q} = d_2^{p + 1, q} \circ d_1^{p, q}$
\end{enumerate}
for all $p, q \in \mathbf{Z}$.
\end{definition}

\noindent
This is just the cochain version of the definition.
It says that each $A^{p, \bullet}$ is a cochain complex
and that each $d_1^{p, \bullet}$ is a morphism of complexes
$A^{p, \bullet} \to A^{p + 1, \bullet}$ such that
$d_1^{p + 1, \bullet} \circ d_1^{p, \bullet} = 0$ as morphisms
of complexes. In other words a double complex can be seen as
a complex of complexes. So in the diagram
$$
\xymatrix{
\ldots &
\ldots &
\ldots &
\ldots \\
\ldots \ar[r] &
A^{p, q + 1} \ar[r]^{d_1^{p, q + 1}} \ar[u] &
A^{p + 1, q + 1} \ar[r] \ar[u] &
\ldots \\
\ldots \ar[r] &
A^{p, q} \ar[r]^{d_1^{p, q}} \ar[u]^{d_2^{p, q}} &
A^{p + 1, q} \ar[r] \ar[u]_{d_2^{p + 1, q}} &
\ldots \\
\ldots &
\ldots \ar[u] &
\ldots \ar[u] &
\ldots
}
$$
any square commutes.
Warning: In the literature one encounters a different definition
where a ``bicomplex'' or  a ``double complex'' has the property that
the squares in the diagram anti-commute.

\begin{example}
\label{example-double-complex-as-tensor-product-of}
Let $\mathcal{A}$, $\mathcal{B}$, $\mathcal{C}$ be additive categories.
Suppose that
$$
\otimes : \mathcal{A} \times \mathcal{B} \longrightarrow \mathcal{C},
\quad
(X, Y) \longmapsto X \otimes Y
$$
is a functor which is bilinear on morphisms, see
Categories, Definition \ref{categories-definition-product-category}
for the definition of $\mathcal{A} \times \mathcal{B}$.
Given complexes $X^\bullet$ of $\mathcal{A}$ and $Y^\bullet$
of $\mathcal{B}$ we obtain a double complex
$$
K^{\bullet, \bullet} = X^\bullet \otimes Y^\bullet
$$
in $\mathcal{C}$. Here the first differential
$K^{p, q} \to K^{p + 1, q}$ is the morphism
$X^p \otimes Y^q \to X^{p + 1} \otimes Y^q$ induced by
the morphism $X^p \to X^{p + 1}$ and the identity on $Y^q$.
Similarly for the second differential.
\end{example}

\begin{definition}
\label{definition-associated-simple-complex}
Let $\mathcal{A}$ be an additive category.
Let $A^{\bullet, \bullet}$ be a double complex.
The {\it associated simple complex}, denoted $sA^\bullet$, also
often called the {\it associated total complex}, denoted
$\text{Tot}(A^{\bullet, \bullet})$, is
given by
$$
sA^n = \text{Tot}^n(A^{\bullet, \bullet}) =
\bigoplus\nolimits_{n = p + q} A^{p, q}
$$
(if it exists) with differential
$$
d_{sA^\bullet}^n = d_{\text{Tot}(A^{\bullet, \bullet})}^n =
\sum\nolimits_{n = p + q} (d_1^{p, q} + (-1)^p d_2^{p, q})
$$
\end{definition}

\noindent
If countable direct sums exist in $\mathcal{A}$ or if for each $n$ at most
finitely many $A^{p, n - p}$ are nonzero, then
$\text{Tot}(A^{\bullet, \bullet})$ exists. Note that
the definition is {\it not} symmetric in the indices $(p, q)$.

\begin{remark}
\label{remark-triple-complex}
Let $\mathcal{A}$ be an additive category. Let $A^{\bullet, \bullet, \bullet}$
be a triple complex. The associated total complex is the complex with
terms
$$
\text{Tot}^n(A^{\bullet, \bullet, \bullet}) =
\bigoplus\nolimits_{p + q + r = n} A^{p, q, r}
$$
and differential
$$
d^n_{\text{Tot}(A^{\bullet, \bullet, \bullet})} =
\sum\nolimits_{p + q + r = n}
d_1^{p, q, r} + (-1)^pd_2^{p, q, r} + (-1)^{p + q}d_3^{p, q, r}
$$
With this definition a simple calculation shows that the associated total
complex is equal to
$$
\text{Tot}(A^{\bullet, \bullet, \bullet}) =
\text{Tot}(\text{Tot}_{12}(A^{\bullet, \bullet, \bullet})) =
\text{Tot}(\text{Tot}_{23}(A^{\bullet, \bullet, \bullet}))
$$
In other words, we can either first combine the first two of the variables
and then combine sum of those with the last, or we can first combine the
last two variables and then combine the first with the sum of the last two.
\end{remark}

\begin{remark}
\label{remark-shift-double-complex}
Let $\mathcal{A}$ be an additive category. Let $A^{\bullet, \bullet}$
be a double complex with differentials $d_1^{p, q}$ and $d_2^{p, q}$.
Denote $A^{\bullet, \bullet}[a, b]$ the double complex with
$$
(A^{\bullet, \bullet}[a, b])^{p, q} = A^{p + a, q + b}
$$
and differentials
$$
d_{A^{\bullet, \bullet}[a, b], 1}^{p, q} = (-1)^a d_1^{p + a, q + b}
\quad\text{and}\quad
d_{A^{\bullet, \bullet}[a, b], 2}^{p, q} = (-1)^b d_2^{p + a, q + b}
$$
In this situation there is a well defined isomorphism
$$
\gamma :
\text{Tot}(A^{\bullet, \bullet})[a + b]
\longrightarrow
\text{Tot}(A^{\bullet, \bullet}[a, b])
$$
which in degree $n$ is given by the map
$$
\xymatrix{
(\text{Tot}(A^{\bullet, \bullet})[a + b])^n =
\bigoplus_{p + q = n + a + b} A^{p, q}
\ar[d]^{\epsilon(p, q, a, b)\text{id}_{A^{p, q}}} \\
\text{Tot}(A^{\bullet, \bullet}[a, b])^n =
\bigoplus_{p' + q' = n} A^{p' + a, q' + b}
}
$$
for some sign $\epsilon(p, q, a, b)$. Of course the summand $A^{p, q}$
maps to the summand $A^{p' + a, q' + b}$ when $p = p' + a$ and $q = q' + b$.
To figure out the conditions on these signs observe that on the source we have
$$
d|_{A^{p, q}} = (-1)^{a + b}\left(d_1^{p, q} + (-1)^pd_2^{p, q}\right)
$$
whereas on the target we have
$$
d|_{A^{p' + a, q' + b}} =
(-1)^ad_1^{p' + a, q' + b} + (-1)^{p'}(-1)^bd_2^{p' + a, q' + b}
$$
Thus our constraints are that
$$
(-1)^a \epsilon(p, q, a, b) = \epsilon(p + 1, q, a, b)(-1)^{a + b}
\Leftrightarrow
\epsilon(p + 1, q, a, b) = (-1)^b \epsilon(p, q, a, b)
$$
and
$$
(-1)^{p' + b}\epsilon(p, q, a, b) =
\epsilon(p, q + 1, a, b) (-1)^{a + b + p}
\Leftrightarrow
\epsilon(p, q, a, b) = \epsilon(p, q + 1, a, b)
$$
Thus we choose $\epsilon(p, q, a, b) = (-1)^{pb}$.
\end{remark}

\begin{remark}
\label{remark-double-complex-complex-of-complexes-first}
Let $\mathcal{A}$ be an additive category with countable direct sums.
Let $\text{DoubleComp}(\mathcal{A})$ denote the category of double complexes.
We can consider an object $A^{\bullet, \bullet}$ of
$\text{DoubleComp}(\mathcal{A})$ as a complex of complexes
as follows
$$
\ldots \to A^{\bullet, -1} \to A^{\bullet, 0} \to A^{\bullet, 1} \to \ldots
$$
For the variant where we switch the role of the indices, see
Remark \ref{remark-double-complex-complex-of-complexes-second}.
In this remark we show that taking the associated total complex
is compatible with all the structures on complexes we have studied
in the chapter so far.

\medskip\noindent
First, observe that the shift functor on double complexes viewed
as complexes of complexes in the manner given above is the functor
$[0, 1]$ defined in Remark \ref{remark-shift-double-complex}.
By Remark \ref{remark-shift-double-complex} the functor
$$
\text{Tot} : \text{DoubleComp}(\mathcal{A}) \to \text{Comp}(\mathcal{A})
$$
is compatible with shift functors, in the sense that we have a functorial
isomorphism $\gamma : \text{Tot}(A^{\bullet, \bullet})[1] \to
\text{Tot}(A^{\bullet, \bullet}[0, 1])$.

\medskip\noindent
Second, if
$$
f, g : A^{\bullet, \bullet} \to B^{\bullet, \bullet}
$$
are homotopic when $f$ and $g$ are viewed as morphisms of complexes
of complexes in the manner given above, then
$$
\text{Tot}(f), \text{Tot}(g) :
\text{Tot}(A^{\bullet, \bullet}) \to \text{Tot}(B^{\bullet, \bullet})
$$
are homotopic maps of complexes. Indeed, let $h = (h^q)$
be a homotopy between $f$ and $g$. If we denote
$h^{p, q} : A^{p, q} \to B^{p, q - 1}$ the component in degree $p$ of $h^q$,
then this means that
$$
f^{p, q} - g^{p, q} = d_2^{p, q - 1} \circ h^{p, q} +
h^{p, q + 1} \circ d_2^{p, q}
$$
The fact that $h^q : A^{\bullet, q} \to B^{\bullet, q - 1}$ is a map of
complexes means that
$$
d_1^{p, q - 1} \circ h^{p, q} = h^{p + 1, q} \circ d_1^{p, q}
$$
Let us define $h' = ((h')^n)$ the homotopy given by the maps
$(h')^n : \text{Tot}^n(A^{\bullet, \bullet}) \to
\text{Tot}^{n - 1}(B^{\bullet, \bullet})$
using $(-1)^ph^{p, q}$ on the summand $A^{p, q}$ for $p + q = n$.
Then we see that
$$
d_{\text{Tot}(B^{\bullet, \bullet})} \circ h' +
h' \circ d_{\text{Tot}(A^{\bullet, \bullet})}
$$
restricted to the summand $A^{p, q}$ is equal to
$$
d_1^{p, q - 1} \circ (-1)^p h^{p, q} +
(-1)^p d_2^{p, q - 1} \circ (-1)^p h^{p, q} +
(-1)^{p + 1} h^{p + 1, q} \circ d_1^{p, q} +
(-1)^p h^{p, q + 1} \circ (-1)^p d_2^{p, q}
$$
which evaluates to $f^{p, q} - g^{p, q}$ by the equations given above.
This proves the second compatibility.

\medskip\noindent
Third, suppose that in the paragraph above we have $f = g$.
Then the assignment $h \leadsto h'$ above is compatible with
the identification of Lemma \ref{lemma-homotopy-shift-cochain}.
More precisely, if we view $h$ as a morphism of complexes
of complexes $A^{\bullet, \bullet} \to B^{\bullet, \bullet}[0, -1]$
via this lemma then
$$
\text{Tot}(A^{\bullet, \bullet})
\xrightarrow{\text{Tot}(h)}
\text{Tot}(B^{\bullet, \bullet}[0, -1])
\xrightarrow{\gamma^{-1}}
\text{Tot}(B^{\bullet, \bullet})[-1]
$$
is equal to $h'$ viewed as a morphism of complexes via the lemma.
Here $\gamma$ is the identification of
Remark \ref{remark-shift-double-complex}.
The verification of this third point is immediate.

\medskip\noindent
Fourth, let
$$
0 \to A^{\bullet, \bullet} \to B^{\bullet, \bullet} \to
C^{\bullet, \bullet} \to 0
$$
be a complex of double complexes and suppose we are given splittings
$s^q : C^{\bullet, q} \to B^{\bullet, q}$ and
$\pi^q : B^{\bullet, q} \to A^{\bullet, q}$
of this as in Lemma \ref{lemma-ses-termwise-split-cochain}
when we view double complexes
as complexes of complexes in the manner given above.
This on the one hand produces a map
$$
\delta : C^{\bullet, \bullet} \longrightarrow A^{\bullet, \bullet}[0, 1]
$$
by the procedure in Lemma \ref{lemma-ses-termwise-split-cochain}.
On the other hand taking $\text{Tot}$ we obtain a complex
$$
0 \to \text{Tot}(A^{\bullet, \bullet}) \to
\text{Tot}(B^{\bullet, \bullet}) \to
\text{Tot}(C^{\bullet, \bullet}) \to 0
$$
which is termwise split (see below) and hence comes with a morphism
$$
\delta' :
\text{Tot}(C^{\bullet, \bullet})
\longrightarrow
\text{Tot}(A^{\bullet, \bullet})[1]
$$
well defined up to homotopy by Lemmas \ref{lemma-ses-termwise-split-cochain}
and \ref{lemma-ses-termwise-split-homotopy-cochain}. Claim:
these maps agree in the sense that
$$
\text{Tot}(C^{\bullet, \bullet})
\xrightarrow{\text{Tot}(\delta)}
\text{Tot}(A^{\bullet, \bullet}[0, 1]) \xrightarrow{\gamma^{-1}}
\text{Tot}(A^{\bullet, \bullet})[1]
$$
is equal to $\delta'$ where $\gamma$ is as in
Remark \ref{remark-shift-double-complex}. To see this denote
$s^{p, q} : C^{p, q} \to B^{\bullet, q}$ and
$\pi^{p, q} : B^{p, q} \to A^{p, q}$ the components of $s^q$ and $\pi^q$.
As splittings
$(s')^n : \text{Tot}^n(C^{\bullet, \bullet}) \to
\text{Tot}^n(B^{\bullet, \bullet})$
and
$(\pi')^n : \text{Tot}^n(B^{\bullet, \bullet}) \to
\text{Tot}^n(A^{\bullet, \bullet})$
we use the maps whose components are $s^{p, q}$ and $\pi^{p, q}$
for $p + q = n$. We recall that
$$
(\delta')^n =
(\pi')^{n + 1} \circ d_{\text{Tot}(B^{\bullet, \bullet})}^n \circ (s')^n :
\text{Tot}^n(C^{\bullet, \bullet}) \to
\text{Tot}^{n + 1}(A^{\bullet, \bullet})
$$
The restriction of this to the summand $C^{p, q}$ is equal to
$$
\pi^{p + 1, q} \circ
d_1^{p, q} \circ
s^{p, q} +
\pi^{p, q + 1} \circ
(-1)^p d_2^{p, q} \circ
s^{p, q} =
\pi^{p, q + 1} \circ
(-1)^p d_2^{p, q} \circ
s^{p, q}
$$
The equality holds because $s^q$ is a morphism of complexes (with $d_1$
as differential) and because $\pi^{p + 1, q} \circ s^{p + 1, q} = 0$
as $s$ and $\pi$ correspond to a direct sum decomposition of $B$
in every bidegree. On the other hand, for $\delta$ we have
$$
\delta^q =  \pi^q \circ d_2 \circ s^q :
C^{\bullet, q} \to A^{\bullet, q + 1}
$$
whose restriction to the summand $C^{p, q}$ is equal to
$\pi^{p, q + 1} \circ d_2^{p, q} \circ s^{p, q}$.
The difference in signs is exactly canceled out by the sign
of $(-1)^p$ in the isomorphism $\gamma$ and the fourth claim is proven.
\end{remark}

\begin{remark}
\label{remark-double-complex-complex-of-complexes-second}
Let $\mathcal{A}$ be an additive category with countable direct sums.
Let $\text{DoubleComp}(\mathcal{A})$ denote the category of double complexes.
We can consider an object $A^{\bullet, \bullet}$ of
$\text{DoubleComp}(\mathcal{A})$ as a complex of complexes
as follows
$$
\ldots \to A^{-1, \bullet} \to A^{0, \bullet} \to A^{1, \bullet} \to \ldots
$$
For the variant where we switch the role of the indices, see
Remark \ref{remark-double-complex-complex-of-complexes-first}.
In this remark we show that taking the associated total complex
is compatible with all the structures on complexes we have studied
in the chapter so far.

\medskip\noindent
First, observe that the shift functor on double complexes viewed
as complexes of complexes in the manner given above is the functor
$[1, 0]$ defined in Remark \ref{remark-shift-double-complex}.
By Remark \ref{remark-shift-double-complex} the functor
$$
\text{Tot} : \text{DoubleComp}(\mathcal{A}) \to \text{Comp}(\mathcal{A})
$$
is compatible with shift functors, in the sense that we have a functorial
isomorphism $\gamma : \text{Tot}(A^{\bullet, \bullet})[1] \to
\text{Tot}(A^{\bullet, \bullet}[1, 0])$.

\medskip\noindent
Second, if
$$
f, g : A^{\bullet, \bullet} \to B^{\bullet, \bullet}
$$
are homotopic when $f$ and $g$ are viewed as morphisms of complexes
of complexes in the manner given above, then
$$
\text{Tot}(f), \text{Tot}(g) :
\text{Tot}(A^{\bullet, \bullet}) \to \text{Tot}(B^{\bullet, \bullet})
$$
are homotopic maps of complexes. Indeed, let $h = (h^p)$
be a homotopy between $f$ and $g$. If we denote
$h^{p, q} : A^{p, q} \to B^{p - 1, q}$ the component in degree $p$ of $h^q$,
then this means that
$$
f^{p, q} - g^{p, q} = d_1^{p - 1, q} \circ h^{p, q} +
h^{p + 1, q} \circ d_1^{p, q}
$$
The fact that $h^p : A^{p, \bullet} \to B^{p - 1, \bullet}$ is a map of
complexes means that
$$
d_2^{p - 1, q} \circ h^{p, q} = h^{p, q + 1} \circ d_2^{p, q}
$$
Let us define $h' = ((h')^n)$ the homotopy given by the maps
$(h')^n : \text{Tot}^n(A^{\bullet, \bullet}) \to
\text{Tot}^{n - 1}(B^{\bullet, \bullet})$
using $h^{p, q}$ on the summand $A^{p, q}$ for $p + q = n$.
Then we see that
$$
d_{\text{Tot}(B^{\bullet, \bullet})} \circ h' +
h' \circ d_{\text{Tot}(A^{\bullet, \bullet})}
$$
restricted to the summand $A^{p, q}$ is equal to
$$
d_1^{p - 1, q} \circ h^{p, q} +
(-1)^{p - 1} d_2^{p - 1, q} \circ h^{p, q} +
h^{p + 1, q} \circ d_1^{p, q} +
h^{p, q + 1} \circ (-1)^p d_2^{p, q}
$$
which evaluates to $f^{p, q} - g^{p, q}$ by the equations given above.
This proves the second compatibility.

\medskip\noindent
Third, suppose that in the paragraph above we have $f = g$.
Then the assignment $h \leadsto h'$ above is compatible with
the identification of Lemma \ref{lemma-homotopy-shift-cochain}.
More precisely, if we view $h$ as a morphism of complexes
of complexes $A^{\bullet, \bullet} \to B^{\bullet, \bullet}[-1, 0]$
via this lemma then
$$
\text{Tot}(A^{\bullet, \bullet})
\xrightarrow{\text{Tot}(h)}
\text{Tot}(B^{\bullet, \bullet}[-1, 0])
\xrightarrow{\gamma^{-1}}
\text{Tot}(B^{\bullet, \bullet})[-1]
$$
is equal to $h'$ viewed as a morphism of complexes via the lemma.
Here $\gamma$ is the identification of
Remark \ref{remark-shift-double-complex}.
The verification of this third point is immediate.

\medskip\noindent
Fourth, let
$$
0 \to A^{\bullet, \bullet} \to B^{\bullet, \bullet} \to
C^{\bullet, \bullet} \to 0
$$
be a complex of double complexes and suppose we are given splittings
$s^p : C^{p, \bullet} \to B^{p, \bullet}$ and
$\pi^p : B^{p, \bullet} \to A^{p, \bullet}$
of this as in Lemma \ref{lemma-ses-termwise-split-cochain}
when we view double complexes
as complexes of complexes in the manner given above.
This on the one hand produces a map
$$
\delta : C^{\bullet, \bullet} \longrightarrow A^{\bullet, \bullet}[0, 1]
$$
by the procedure in Lemma \ref{lemma-ses-termwise-split-cochain}.
On the other hand taking $\text{Tot}$ we obtain a complex
$$
0 \to \text{Tot}(A^{\bullet, \bullet}) \to
\text{Tot}(B^{\bullet, \bullet}) \to
\text{Tot}(C^{\bullet, \bullet}) \to 0
$$
which is termwise split (see below) and hence comes with a morphism
$$
\delta' :
\text{Tot}(C^{\bullet, \bullet})
\longrightarrow
\text{Tot}(A^{\bullet, \bullet})[1]
$$
well defined up to homotopy by Lemmas \ref{lemma-ses-termwise-split-cochain}
and \ref{lemma-ses-termwise-split-homotopy-cochain}. Claim:
these maps agree in the sense that
$$
\text{Tot}(C^{\bullet, \bullet})
\xrightarrow{\text{Tot}(\delta)}
\text{Tot}(A^{\bullet, \bullet}[1, 0]) \xrightarrow{\gamma^{-1}}
\text{Tot}(A^{\bullet, \bullet})[1]
$$
is equal to $\delta'$ where $\gamma$ is as in
Remark \ref{remark-shift-double-complex}. To see this denote
$s^{p, q} : C^{p, q} \to B^{\bullet, q}$ and
$\pi^{p, q} : B^{p, q} \to A^{p, q}$ the components of $s^q$ and $\pi^q$.
As splittings
$(s')^n : \text{Tot}^n(C^{\bullet, \bullet}) \to
\text{Tot}^n(B^{\bullet, \bullet})$
and
$(\pi')^n : \text{Tot}^n(B^{\bullet, \bullet}) \to
\text{Tot}^n(A^{\bullet, \bullet})$
we use the maps whose components are $s^{p, q}$ and $\pi^{p, q}$
for $p + q = n$. We recall that
$$
(\delta')^n =
(\pi')^{n + 1} \circ d_{\text{Tot}(B^{\bullet, \bullet})}^n \circ (s')^n :
\text{Tot}^n(C^{\bullet, \bullet}) \to
\text{Tot}^{n + 1}(A^{\bullet, \bullet})
$$
The restriction of this to the summand $C^{p, q}$ is equal to
$$
\pi^{p + 1, q} \circ
d_1^{p, q} \circ
s^{p, q} +
\pi^{p, q + 1} \circ
(-1)^p d_2^{p, q} \circ
s^{p, q} =
\pi^{p + 1, q} \circ
d_1^{p, q} \circ
s^{p, q}
$$
The equality holds because $s^p$ is a morphism of complexes (with $d_2$
as differential) and because $\pi^{p, q + 1} \circ s^{p, q + 1} = 0$
as $s$ and $\pi$ correspond to a direct sum decomposition of $B$
in every bidegree. On the other hand, for $\delta$ we have
$$
\delta^p =  \pi^p \circ d_1 \circ s^p :
C^{p, \bullet} \to A^{p + 1, \bullet}
$$
whose restriction to the summand $C^{p, q}$ is equal to
$\pi^{p + 1, q} \circ d_1^{p, q} \circ s^{p, q}$.
Thus we get the same as before which matches with the fact that
the isomorphism
$\gamma : \text{Tot}(A^{\bullet, \bullet})[1] \to
\text{Tot}(A^{\bullet, \bullet}[1, 0])$
is defined without the intervention of signs.
\end{remark}







